\documentclass[11pt,a4paper]{article}
\usepackage{amsmath,amssymb,amsthm,amsfonts,bm}
\usepackage{geometry}
\geometry{margin=1in}

% Theorem environments
\theoremstyle{plain}
\newtheorem{axiom}{Axiom}[section]
\newtheorem{theorem}{Theorem}[section]
\newtheorem{conjecture}{Conjecture}[section]

\title{Semantic Validation Gates: A Formal Axiomatic Framework for Semantic Verification}
\author{Luoqing ZHAO\\Independent Researcher}
\date{June 2026}

\begin{document}
	
	\maketitle
	
	\noindent\textbf{Status:} Work in Progress – Preliminary Draft\\
	\textbf{SSRN Classification:} Computer Science / Artificial Intelligence / Logic \& Verification
	
	\vspace{1em}
	
	\begin{abstract}
As large language models and autonomous agents generate increasingly complex reasoning chains, explanations, and decisions, there is a growing need for a systematic, mathematically grounded way to verify their outputs. This paper introduces a formal axiomatic framework of five validation gates that transforms qualitative judgments of correctness, consistency, and alignment into geometrically well‑defined measures of deviation. The framework guarantees mutual exclusivity and completeness of error classification, providing both a rigorous theoretical foundation and a clear path toward empirical calibration. It is designed to support reliable reasoning engineering, loop stability control, and systematic testing harness construction for automated AI systems.

	\end{abstract}
	
	\section{Introduction}
	
	Large language models (LLMs) and autonomous agents have demonstrated remarkable capabilities across reasoning, generation, and decision-making tasks. However, their widespread adoption reveals two fundamental limitations: \textit{computational inefficiency} and \textit{output inconsistency}. Operating primarily at the token level, these systems generate responses sequentially, often consuming significant compute and memory resources even for routine tasks. More critically, their outputs are inherently probabilistic. They can appear fluent and structurally correct while containing factual errors, logical contradictions, safety risks, or misalignment with the original intent. This lack of systematic verification creates uncertainty, increases the need for repeated regeneration, and further drives up token consumption and operational costs.
	
	In traditional computing and digital systems, reliability is enforced through well-defined logic gates and deterministic checks, providing clear, repeatable rules to validate every operation. No equivalent standardized framework exists for the probabilistic, semantic outputs of modern artificial intelligence. Current evaluation methods are often ad-hoc, task-specific, or limited to a single dimension—such as factuality or toxicity—without a unified structure that spans from basic syntax to final goal alignment.
	
	To address this gap, we propose \textbf{Semantic Validation Gates}: a formal axiomatic framework that extends the concept of structured verification to the semantic space of AI outputs. By defining a sequence of mathematically grounded validation operators, we transform qualitative judgments into measurable deviations. This enables systematic verification across five complementary dimensions: format compliance, factual consistency, logical coherence, safety alignment, and intent fulfillment. This layered, lexicographic approach provides a clear, standardized path to identify exactly where and why an output fails, thereby reducing unnecessary regeneration and improving both efficiency and accuracy.
	
	This paper presents the complete mathematical foundation of the framework, including its core axioms, operator definitions, a rigorous completeness theorem, and formal conjectures guiding future topological and metric development.
	
	\section{Axiomatic Foundation}
	
	We begin with three core axioms defining the mathematical structure of the framework.
	
	\begin{axiom}[Geometric Structure]
		The semantic space is modeled as a $d$-dimensional complete Riemannian manifold $(\mathcal{M}, g)$ equipped with the Fisher information metric. Let $\mathcal{P}(\mathcal{M})$ denote the space of Borel probability measures on $\mathcal{M}$, endowed with the $2$-Wasserstein metric $W_2$, which quantifies the dissimilarity between distributions. Any model output is represented as a measure $\mu \in \mathcal{P}(\mathcal{M})$.
	\end{axiom}
	
	\begin{axiom}[Fiber Bundle Structure]
		There exists a fiber bundle $E \xrightarrow{\pi} B$ over a base space $B$, allowing local trivialization: for any open semantic patch $U_\alpha \subset B$, there is a homeomorphism
		\[
		\Phi_\alpha: \pi^{-1}(U_\alpha) \to U_\alpha \times F,
		\]
		where $F$ is the standard fiber. The mapping $\Phi$ is defined via Karcher barycentric projection, ensuring consistency across overlapping regions.
	\end{axiom}
	
	\begin{axiom}[Normalized Deviation]
		For each validation dimension $i = 1, \dots, 5$, we define a relative deviation operator:
		\[
		\rho_i(\mu) = \frac{f_i(\mu)}{\tau_i},
		\]
		where $f_i: \mathcal{P}(\mathcal{M}) \to [0, +\infty)$ is the raw deviation measure, and $\tau_i > 0$ is a tolerance threshold calibrated using spectral gaps or empirical distributions. A dimension is considered valid if and only if $\rho_i(\mu) \le 1$.
	\end{axiom}
	
	\section{The Five Validation Operators}
	
	We define the set of verification operators $\mathcal{F} = \{f_1, f_2, f_3, f_4, f_5\}$ as follows.
	
	\noindent\textbf{$f_1$: Format Gate}
	\[
	f_1(\mu) = W_2\left(\mu, \operatorname{Proj}_{\mathcal{P}_{\text{form}}}(\mu)\right),
	\]
	where $\mathcal{P}_{\text{form}} \subset \mathcal{P}(\mathcal{M})$ is a closed, convex subset containing all distributions satisfying syntactic and structural constraints. This measures the distance from the nearest valid structure.
	
	\vspace{0.5em}
	\noindent\textbf{$f_2$: Fact Gate}
	\[
	f_2(\mu) = W_2\left(\mu, \nu_{\text{fact}}\right),
	\]
	where $\nu_{\text{fact}}$ is a reference measure derived from a verified knowledge base, smoothed via kernel density estimation (KDE). This measures deviation from established facts.
	
	\vspace{0.5em}
	\noindent\textbf{$f_3$: Logic Gate}
	Let $s_\mu = \Phi(\mu)$ be the local section representation of $\mu$, and let $o(s_\mu) \in \check{H}^1(\mathcal{U}, \mathcal{S})$ denote the Čech cohomology obstruction class encoding internal inconsistencies. Then:
	\[
	f_3(\mu) = \left\| \mathcal{H}\left(o(s_\mu)\right) \right\|_{\text{harm}},
	\]
	where $\mathcal{H}$ is the Hodge–Laplace operator, and $\|\cdot\|_{\text{harm}}$ is the norm on harmonic forms. This quantifies the severity of logical contradictions.
	
	\vspace{0.5em}
	\noindent\textbf{$f_4$: Alignment Gate}
	\[
	f_4(\mu) = W_2\left(\mu, \operatorname{Proj}_{\mathcal{P}_{\text{safe}}}(\mu)\right),
	\]
	where $\mathcal{P}_{\text{safe}} \subset \mathcal{P}(\mathcal{M})$ is a compact, convex set representing safe and aligned behaviors. Smaller values indicate closer adherence to safety guidelines, consistent with the validity condition $\rho_4(\mu) \le 1$.
	
	\vspace{0.5em}
	\noindent\textbf{$f_5$: Intent Gate}
	\[
	f_5(\mu) = \left\| \psi(\mu) - \psi(\nu_{\text{intent}}) \right\|_{\mathcal{H}_{\text{RKHS}}},
	\]
	where $\psi: \mathcal{P}(\mathcal{M}) \to \mathcal{H}_{\text{RKHS}}$ is an isometric embedding into a reproducing kernel Hilbert space, and $\nu_{\text{intent}}$ is the target distribution representing the task objective. This measures alignment with the original goal.
	
	\section{Validity and Error Partitioning}
	
	\noindent\textbf{Valid Space}\\
	The set of fully valid outputs is the intersection of all acceptable deviations:
	\[
	\mathcal{V}_{\text{valid}} = \bigcap_{i=1}^5 \left\{ \mu \in \mathcal{P}(\mathcal{M}) \,\bigg|\, \rho_i(\mu) \le 1 \right\}.
	\]
	
	\vspace{0.5em}
	\noindent\textbf{Priority Mapping}\\
	Define the classification function:
	\[
	\Pi(\mu) = \min\left\{ i \in \{1,\dots,5\} \,\bigg|\, \rho_i(\mu) > 1 \right\},
	\]
	with $\Pi(\mu) = 0$ if all $\rho_i(\mu) \le 1$. This enforces a lexicographic evaluation order: \textbf{Format $\to$ Fact $\to$ Logic $\to$ Alignment $\to$ Intent}.
	
	\vspace{0.5em}
	\noindent\textbf{Error Subsets}\\
	For each $i = 1, \dots, 5$, the distinct error regions are:
	\[
	\mathcal{E}_i = \left\{ \mu \in \mathcal{P}(\mathcal{M}) \,\bigg|\, \rho_i(\mu) > 1 \quad \text{and} \quad \forall j < i,\ \rho_j(\mu) \le 1 \right\}.
	\]
	
	\section{Completeness and Consistency Properties}
	
	\subsection{Proven Result: Completeness and Mutual Exclusivity}
	
	\begin{theorem}[Completeness Theorem]
		The space of all possible model outputs decomposes into disjoint valid and error regions:
		\[
		\mathcal{P}(\mathcal{M}) = \mathcal{V}_{\text{valid}} \sqcup \bigcup_{i=1}^5 \mathcal{E}_i,
		\]
		where $\sqcup$ denotes disjoint union. Equivalently:
		\[
		\mathcal{E} := \mathcal{P}(\mathcal{M}) \setminus \mathcal{V}_{\text{valid}} = \bigcup_{i=1}^5 \mathcal{E}_i.
		\]
	\end{theorem}
	
	\begin{proof}
		\begin{enumerate}
			\item \textbf{$\bigcup_{i=1}^5 \mathcal{E}_i \subseteq \mathcal{E}$:}\\
			If $\mu \in \mathcal{E}_i$, then by definition $\rho_i(\mu) > 1$, which implies $\mu \notin \mathcal{V}_{\text{valid}}$. Hence, $\mu \in \mathcal{E}$.
			
			\item \textbf{$\mathcal{E} \subseteq \bigcup_{i=1}^5 \mathcal{E}_i$:}\\
			If $\mu \in \mathcal{E}$, then the set of failed gates $K(\mu) = \{i : \rho_i(\mu) > 1\}$ is non-empty. By the well-ordering principle of natural numbers, there exists a unique minimum index $k_{\min} = \min K(\mu)$. It follows that $\rho_{k_{\min}}(\mu) > 1$ and $\rho_j(\mu) \le 1$ for all $j < k_{\min}$. Therefore, $\mu \in \mathcal{E}_{k_{\min}}$.
			
			\item \textbf{Mutual Exclusivity:}\\
			For $i \neq j$, assume without loss of generality that $i < j$. Any output $\mu \in \mathcal{E}_i$ satisfies $\rho_i(\mu) > 1$, so it cannot satisfy $\rho_i(\mu) \le 1$, which is required for membership in $\mathcal{E}_j$. Thus, $\mathcal{E}_i \cap \mathcal{E}_j = \emptyset$.
		\end{enumerate}
		This completes the proof.
	\end{proof}
	
	\subsection{Analytical Properties to Be Proved (Work in Progress)}
	
	To fully establish the mathematical consistency, topological stability, and empirical robustness of the framework, the following statements are posed as formal conjectures to be derived and proved in subsequent work:
	
	\begin{conjecture}[Boundary Regularity]
		Let $\lambda$ be the reference measure on the Wasserstein space $\mathcal{P}(\mathcal{M})$. The boundary of the valid region satisfies:
		\[
		\dim_H(\partial\mathcal{V}_{\text{valid}}) \le d-1 \quad \text{and} \quad \lambda(\partial\mathcal{V}_{\text{valid}}) = 0,
		\]
		where $\dim_H$ denotes the Hausdorff dimension.
	\end{conjecture}
	\noindent\textit{Implication:} The boundary does not affect the almost-everywhere well-posedness of the classification mapping $\Pi(\mu)$.
	
	\begin{conjecture}[Metric Consistency and Bounds]
		Assume $\mathcal{M}$ satisfies a Bakry–Émery Ricci curvature lower bound $\operatorname{Ric}_\infty \ge \kappa > 0$. Then all distance functionals $W_2$, $\|\cdot\|_{\text{harm}}$, and $\|\cdot\|_{\mathcal{H}_{\text{RKHS}}}$ are globally Lipschitz equivalent. Consequently, each deviation operator $f_i$ is Lipschitz continuous with constants uniformly bounded over $\mathcal{P}(\mathcal{M})$.
	\end{conjecture}
	\noindent\textit{Implication:} Small changes in the input distribution lead to controlled changes in measured deviation, ensuring stability.
	
	\begin{conjecture}[Threshold Well-Posedness via EVT]
		For any significance level $\alpha \in (0,1)$, there exists a unique tolerance vector $\boldsymbol{\tau} = (\tau_1, \dots, \tau_5)$ derived from the Fréchet distribution of extreme deviations over valid reference samples, such that:
		\[
		\mathbb{P}\left( \rho_i(\mu) > 1 \,\big|\, \mu \text{ valid} \right) \le \alpha.
		\]
	\end{conjecture}
	\noindent\textit{Implication:} Thresholds are statistically grounded and bound the probability of false rejection.
	
	\begin{conjecture}[Continuity Under Perturbation]
		The classification mapping $\Pi: \mathcal{P}(\mathcal{M}) \to \{0,1,\dots,5\}$ is locally constant on the interior of each partition $\mathcal{V}_{\text{valid}}$ and $\mathcal{E}_i$. That is, for any $\mu$ in the interior of a region, there exists $\varepsilon > 0$ such that:
		\[
		W_2(\mu, \mu') < \varepsilon \quad \Rightarrow \quad \Pi(\mu') = \Pi(\mu).
		\]
	\end{conjecture}
	\noindent\textit{Implication:} Classification is robust against minor semantic noise and sampling variation.
	
	\section{Current Status \& Next Steps}
	
	This draft presents the complete axiomatic structure, operator definitions, and core completeness theorem. Ongoing work includes formal proofs of the conjectures above, development of the calibration methodology, and empirical validation on reasoning benchmarks such as GSM8K, MATH, and TruthfulQA.
	
	\vspace{1em}
	\noindent\textbf{Disclaimer:} This is a preliminary working paper. Comments, suggestions, and collaboration inquiries are welcome. All rights reserved.
	
\end{document}
